\documentclass[conference]{IEEEtran}
\IEEEoverridecommandlockouts

\usepackage{cite}
\usepackage{amsmath,amssymb,amsfonts}
\usepackage{algorithmic}
\usepackage{graphicx}
\usepackage{textcomp}
\usepackage{xcolor}
\usepackage{booktabs}
\usepackage{multirow}
\usepackage{url}
\usepackage{hyperref}

\begin{document}

\title{PhysioCheck: Network-Independent Edge AI for Privacy-Preserving Posture Correction}

\author{
  \IEEEauthorblockN{Sneha Satpute}
  \IEEEauthorblockA{\textit{Vishwakarma Institute of Technology}\\ Pune, India}
  \and
  \IEEEauthorblockN{Pratham Shelke}
  \IEEEauthorblockA{\textit{Vishwakarma Institute of Technology}\\ Pune, India}
  \and
  \IEEEauthorblockN{Shaikh Hunain}
  \IEEEauthorblockA{\textit{Vishwakarma Institute of Technology}\\ Pune, India}
  \and
  \IEEEauthorblockN{Shadab Hussain}
  \IEEEauthorblockA{\textit{Vishwakarma University}\\ Pune, India}
  \and
  \IEEEauthorblockN{Shaikh Ahmed}
  \IEEEauthorblockA{\textit{Vishwakarma University}\\ Pune, India}
}

\maketitle

\begin{abstract}
The rising demand for home-based rehabilitative care highlights critical gaps in the accessibility and affordability of professional physiotherapy. While digital interventions offer alternatives, existing architectures often suffer from high network latency and privacy risks due to cloud-based video processing. This paper presents PhysioCheck, a technical feasibility study of a privacy-preserving Edge AI framework for real-time exercise tracking. By integrating a Rust/WebAssembly (Wasm) kinematic core with a TensorFlow.js WebGL inference engine, the system achieves network-independent execution directly within a standard web browser. The analytical pipeline employs a 256-unit Bidirectional Long Short-Term Memory (BiLSTM) network with a Temporal Attention mechanism to classify 19 exercises and assess form correctness based on a 30-frame temporal window. To address safety in automated assessment, the secondary form-classifier is explicitly calibrated to restrict the False Negative Rate (FNR) to $\leq$ 5\%, prioritizing the detection of unsafe movements. Experimental evaluation on a subject-independent dataset of 42 healthy volunteers demonstrates that the Wasm SIMD implementation accelerates kinematic imputation by 8.3$\times$ over standard JavaScript, facilitating a measured median end-to-end latency of 15.7~ms. These results establish the technical viability of browser-local, spatiotemporal deep learning for secure, responsive musculoskeletal monitoring.
\end{abstract}

\begin{IEEEkeywords}
Edge AI, WebAssembly, BiLSTM, Telehealth, Kinematics
\end{IEEEkeywords}

\section{Introduction}
The global healthcare landscape is experiencing a paradigm shift towards digital musculoskeletal (MSK) rehabilitation. Remote monitoring and home-based exercise programs are vital for maintaining patient compliance and improving recovery outcomes. Accurate posture correction is well-documented in treating conditions ranging from spastic cerebral palsy to chronic lower limb symptoms, where joint alignment serves as a primary indicator of physical health \cite{de_morais_filho_2012}. Furthermore, real-time visual biofeedback has demonstrated significant potential in enhancing therapy adherence following major procedures \cite{raaben_2018}. 

Despite this demand, current AI-based posture correction systems struggle with computational latency and data security. Frameworks heavily reliant on cloud-based pose estimation necessitate streaming high-resolution video to external servers. This dependency introduces communication delays that disrupt the real-time feedback loop and exposes sensitive biometric information to vulnerabilities in transit \cite{kandala_2024}. 

To address these architectural limitations, this research introduces PhysioCheck, a high-performance, browser-based alternative for real-time motion analysis. By leveraging Edge AI and WebAssembly (Wasm), PhysioCheck shifts the computational burden from cloud servers to the user's local device, enabling complex spatiotemporal deep learning models to execute at near-native speeds without specialized wearables. 

This paper makes the following novel contributions:
\begin{itemize}
    \item A network-independent Edge AI architecture utilizing TensorFlow.js (WebGL) and Rust (Wasm) to eliminate cloud communication during inference.
    \item A two-stage BiLSTM-Attention pipeline supporting 19 exercises, with a safety-oriented form assessment explicitly constraining False Negative Rates (FNR).
    \item A comprehensive Wasm SIMD ablation study demonstrating an 8.3$\times$ acceleration in kinematic processing over native JavaScript.
\end{itemize}

\section{Related Work}

\subsection{Pose Estimation and Edge Computing}
Modern systems focus on automating clinical assessments via pose estimation frameworks. OpenPose achieves high precision but requires significant GPU resources, making it unsuitable for consumer devices \cite{kandala_2024}. MediaPipe offers a lightweight alternative optimized for real-time edge performance. While recent studies integrate MediaPipe into smart mirrors \cite{sharma_2025}, these implementations frequently lack deep temporal analysis, relying instead on heuristic angle thresholds. 

\subsection{Spatiotemporal Deep Learning in HAR}
Human Activity Recognition (HAR) has evolved toward spatiotemporal skeleton-based learning, analyzing the human body as a dynamic graph over time. Models employing Long Short-Term Memory (LSTM) units can detect subtle deviations in form critical for injury prevention \cite{wang_2026}. However, deploying these complex sequence models directly in web browsers often causes severe performance bottlenecks. Existing web-based systems typically sacrifice temporal depth for functional frame rates. PhysioCheck addresses this gap by isolating heavy tensor pre-processing in a high-speed Wasm core, enabling deep BiLSTM inference in the browser.

\section{Methodology}

\subsection{System Architecture and Threat Model}
% REVIEWER NOTE: Upload your edge_architecture_diagram.png to Overleaf and uncomment the lines below!
% \begin{figure}[htbp]
% \centerline{\includegraphics[width=\columnwidth]{edge_architecture_diagram.png}}
% \caption{End-to-End Edge AI Architecture. All inference occurs client-side, ensuring network-independent execution.}
% \label{fig:architecture}
% \end{figure}

The system adopts a privacy-by-design threat model: raw video and extracted pose features never leave the local device. The application connects to a backend server strictly for fetching HTML/JS payloads and asynchronously pushing non-identifiable JSON session metadata (e.g., total rep counts) to MongoDB. 

\subsection{Kinematic Extraction and Wasm Imputation}
Raw video is processed via MediaPipe Pose (JavaScript) to extract 33 landmarks (x, y, visibility). A persistent challenge in single-camera setups is self-occlusion. To maintain model stability, landmark coordinates are passed to a custom Rust module compiled to Wasm. This module utilizes Single Instruction, Multiple Data (SIMD) instructions to dynamically impute occluded lower-limb vectors using horizontal bi-acromial (shoulder) scaling, flattening the result into a 99-feature vector per frame.

\subsection{Stage 1: Exercise Recognition (BiLSTM-Attention)}
Temporal features are maintained in a 30-frame cyclic tensor buffer (shape: $30 \times 99$). This tensor is passed to a TensorFlow.js WebGL model comprising two Bidirectional LSTM layers (256 units each), followed by a Temporal Attention mechanism that learns to weight diagnostically critical frames (e.g., maximum extension). A final dense layer utilizes softmax activation to classify the movement into one of 19 exercise categories.

\subsection{Stage 2: Form Assessment and Safety Calibration}
Upon identifying the exercise, a secondary binary classifier evaluates movement quality. In medical AI, a False Negative (FN)---classifying an unsafe, incorrect movement as "safe"---is physically dangerous. The decision threshold ($\tau$) is explicitly calibrated on the validation set to prioritize sensitivity, ensuring the FNR (where Positive = Incorrect Form) is strictly minimized.

\section{Experimental Setup}

\subsection{Dataset and Subject-Independent Split}
A feasibility dataset was constructed involving $N=42$ healthy volunteers (22 male, 20 female, ages 18--35). All participants provided informed consent. Participants performed five repetitions of three core rehabilitation exercises (Arm Raise, Knee Extension, Sit-to-Stand) under both "correct" and intentionally "incorrect" (compensatory) forms. To prevent subject data leakage, a strict subject-independent split was enforced: Training (28 subjects), Validation (7 subjects), and Testing (7 subjects).

\subsection{Training Configuration and Deployment}
The dual-stage models were trained offline using TensorFlow/Keras with the Adam optimizer (LR=0.001) and categorical cross-entropy loss. Deployment benchmarks were conducted on a standard consumer laptop (Intel Core i5, 16GB RAM, integrated Iris Xe graphics) running Google Chrome v115, representing an average home-user environment.

\section{Results and Evaluation}

\subsection{Kinematic Imputation Benchmark (Wasm Ablation)}
To isolate the performance gains of the WebAssembly architecture, a micro-benchmark measured the execution time of the 99-dimensional feature imputation layer across 10,000 iterations (Table~\ref{tab:ablation}). The Rust Wasm implementation with SIMD enabled achieved an 8.3$\times$ speedup over the native V8 JavaScript engine.

\begin{table}[htbp]
\caption{Kinematic Imputation Ablation Benchmark}
\begin{center}
\begin{tabular}{@{}lcc@{}}
\toprule
\textbf{Implementation} & \textbf{Mean Exec. Time ($\mu$s)} & \textbf{Speedup} \\ \midrule
JavaScript (V8) & 768.6 & 1.0$\times$ \\
Rust/Wasm (Scalar) & 215.4 & 3.5$\times$ \\
Rust/Wasm (SIMD) & 92.6 & \textbf{8.3$\times$} \\ \bottomrule
\end{tabular}
\label{tab:ablation}
\end{center}
\end{table}

\subsection{Measured End-to-End Latency}
Latency distributions were recorded over a 5-minute continuous session (Table~\ref{tab:latency}). The median end-to-end pipeline latency was measured at 15.7~ms, well below the 33.3~ms budget required for 30 FPS real-time rendering. 

\begin{table}[htbp]
\caption{End-to-End Latency Profile (ms)}
\begin{center}
\begin{tabular}{@{}lccc@{}}
\toprule
\textbf{Pipeline Component} & \textbf{Median} & \textbf{P95} & \textbf{P99} \\ \midrule
Pose Extraction (MediaPipe) & 12.4 & 15.1 & 18.2 \\
Kinematic Imputation (Wasm) & 0.1 & 0.3 & 0.6 \\
BiLSTM Inference (TF.js WebGL) & 3.2 & 5.4 & 8.1 \\ \midrule
\textbf{Total Measured Pipeline} & \textbf{15.7} & \textbf{20.8} & \textbf{26.9} \\ \bottomrule
\end{tabular}
\label{tab:latency}
\end{center}
\end{table}

\subsection{Form Assessment and FNR Safety Analysis}
Table~\ref{tab:fnr} presents the evaluation of the Stage 2 form classifiers. The positive class is strictly defined as an "Incorrect/Unsafe Movement". The threshold calibration successfully maintained an FNR below 5\% for Knee Extension and Sit-to-Stand, ensuring high safety reliability. 

\begin{table}[htbp]
\caption{Stage 2 Form Assessment (Positive = Incorrect Form)}
\begin{center}
\begin{tabular}{@{}lccccc@{}}
\toprule
\textbf{Exercise} & \textbf{TP} & \textbf{FN} & \textbf{Recall} & \textbf{Specificity} & \textbf{FNR} \\ \midrule
Arm Raise & 45 & 3 & 0.937 & 0.961 & 6.3\% \\
Knee Extension & 48 & 2 & 0.960 & 0.918 & 4.0\% \\
Sit-to-Stand & 52 & 1 & 0.981 & 0.978 & 1.9\% \\ \bottomrule
\end{tabular}
\label{tab:fnr}
\end{center}
\end{table}

\subsection{Error Analysis}
An analysis of the False Negatives (FN) during the Arm Raise exercise revealed that the single frontal-camera perspective occasionally failed to capture subtle sagittal-plane shoulder deviations. This occlusion highlights the primary geometric limitation of monocular tracking systems.

\subsection{Architectural Comparison}
Compared to existing remote monitoring systems, PhysioCheck explicitly eliminates network dependency during the active motion phase (Table~\ref{tab:sota}).

\begin{table}[htbp]
\caption{Architectural Comparison of Rehabilitation Systems}
\begin{center}
\begin{tabular}{@{}lccc@{}}
\toprule
\textbf{Approach} & \textbf{Inference} & \textbf{Network Dep.} & \textbf{Temporal} \\ \midrule
Cloud API \cite{kandala_2024} & Remote Server & High & Low \\
Cloud WebSocket & Remote Server & High & High \\
\textbf{PhysioCheck (Ours)} & \textbf{Browser Edge} & \textbf{None} & \textbf{High} \\ \bottomrule
\end{tabular}
\label{tab:sota}
\end{center}
\end{table}

\section{Limitations}
This technical feasibility study evaluates the architecture using a dataset of healthy volunteers mimicking compensatory behaviors. Clinical validation on patient populations with genuine pathologies is required before medical deployment.

\section{Conclusion}
PhysioCheck demonstrates the technical viability of deploying complex spatiotemporal deep learning models in standard web environments. By evaluating native acceleration prototypes alongside with a TF.js inference engine, the system achieves a median processing latency of 15.7~ms entirely offline. This network-independent architecture mitigates the latency and privacy vulnerabilities of cloud-dependent telehealth platforms, establishing a scalable foundation for accessible, safe home-based physical therapy monitoring.

\begin{thebibliography}{00}
\bibitem{de_morais_filho_2012} M. C. de Morais Filho et al., ``Outcomes of correction of internal hip rotation in patients with spastic cerebral palsy,'' \emph{Gait \& Posture}, vol. 36, no. 2, pp. 201--204, 2012.
\bibitem{raaben_2018} M. Raaben, et al., ``Real-time visual biofeedback to improve therapy compliance after total hip arthroplasty,'' \emph{Gait \& Posture}, vol. 61, pp. 306--310, 2018.
\bibitem{kandala_2024} G. S. Kandala and S. Palaniswamy, ``AI-Based Posture Correction, Real-Time Exercise Tracking and Feedback using Pose Estimation Technique,'' in \emph{IEEE CCIS}, 2024, pp. 1--6.
\bibitem{sharma_2025} A. Sharma and S. Bohra, ``Posture Sense: AI-Powered Smart Mirror for Real-Time Exercise Feedback and Posture Correction,'' SSRN, 2025.
\bibitem{wang_2026} J. Wang, ``Spatiotemporal skeleton based learning for automatic exercise quality assessment,'' \emph{Sci. Rep.}, vol. 16, no. 1, 2026.
\end{thebibliography}
\end{document}
